\documentclass[8pt,a4paper,landscape]{extarticle}
\usepackage{cheatsheet-style}   % all styling lives here
\newcommand{\Srn}{\left\{\begin{matrix}r\\ n\end{matrix}\right\}}
\newcommand{\Srk}{\left\{\begin{matrix}r\\ k\end{matrix}\right\}}
\newcommand{\br}[2]{\left[\begin{matrix}#1\\ #2\end{matrix}\right]}
\newcommand{\bc}[2]{\left\{\begin{matrix}#1\\ #2\end{matrix}\right\}}
\newcommand{\Zge}{\mathbb{Z}_{\ge 0}}
\begin{document}
\begin{cscolumns}
%% ======================= CH 1 =======================
\chapterhead*{Ch 1. Basic Counting Principles}

\topic{Pascal}
\begin{theorem}[Pascal's Identity]
\[ \binom{n}{r}=\binom{n-1}{r-1}+\binom{n-1}{r} \]
\end{theorem}

\topic{Euler's $\phi$-function}
Let $n\in\mathbb{N}$, and let
\[ n=p_1^{m_1}p_2^{m_2}\cdots p_k^{m_k} \]
be its prime factorization as stated in the fundamental theorem of arithmetic, with $p_1,\dots,p_k$ distinct primes.
\begin{theorem}[Euler]
\[ \phi(n)=n\left(1-\frac{1}{p_1}\right)\left(1-\frac{1}{p_2}\right)\cdots\left(1-\frac{1}{p_k}\right) \]
\[ =n\prod_{i=1}^{k}\left(1-\frac{1}{p_i}\right). \]
\end{theorem}

%% ======================= CH 2 =======================
\chapterhead*{Ch 2. The Twelvefold Way}

\topic{Rota's Twelvefold Way}
\begin{formula}
\begin{cstable}{>{\raggedright\arraybackslash}p{0.25\linewidth}CCC}
\textbf{Objects / Boxes} & \textbf{No Restriction} & \textbf{Injective} & \textbf{Surjective} \\ \midrule
Obj.\ Dist., Box.\ Dist. & $n^r$ & $P^n_r$ & $n!\Srn$ \\ \midrule
Obj.\ Dist., Box.\ Iden. & $\displaystyle\sum_{k=0}^{n}\Srk$ & $\begin{cases}1 & r\le n\\ 0 & r>n\end{cases}$ & $\Srn$ \\ \midrule
Obj.\ Iden., Box.\ Dist. & $\dbinom{r+n-1}{r}$ & $\dbinom{n}{r}$ & $\dbinom{r-1}{n-1}$ \\ \midrule
Obj.\ Iden., Box.\ Iden. & $\displaystyle\sum_{k=0}^{n}p(r,k)$ & $\begin{cases}1 & r\le n\\ 0 & r>n\end{cases}$ & $p(r,n)$
\end{cstable}
\end{formula}
Here $P^n_r=\frac{n!}{(n-r)!}$ is the number of $r$-permutations of an $n$-set (Chapter 1); $\Srn$, also written $S(r,n)$, is the Stirling number of the second kind; and $p(r,n)$ is the number of partitions of $r$ identical objects into exactly $n$ identical boxes.

\topic{Stirling Numbers of the Second Kind}
The Stirling number of the second kind, denoted by $S(r,n)$ or $\Srn$, is the number of ways to partition a set of $r$ distinct objects into $n$ nonempty, unlabelled subsets.

The Stirling numbers of the second kind satisfy the following conditions:
\begin{formula}
\begin{enumerate}[label=\arabic*.]
\item $S(0,0)=1$, $S(r,0)=0$ for $r>0$, and $S(0,n)=0$ for $n>0$;
\item $S(r,n)=nS(r-1,n)+S(r-1,n-1)$ for $r,n\ge 1$.
\end{enumerate}
\end{formula}
Unlike the binomial coefficients, the Stirling numbers have no product formula. There is, however, an explicit formula as an alternating sum, which comes from the principle of inclusion and exclusion.
\begin{formula}
For all $r,n\ge 0$,
\[ S(r,n)=\frac{1}{n!}\sum_{j=0}^{n}(-1)^{n-j}\binom{n}{j}j^r. \]
\end{formula}
\begin{formula}
\[ \br{n+1}{k}=n\br{n}{k}+\br{n}{k-1},\quad \bc{n+1}{k}=k\bc{n}{k}+\bc{n}{k-1}. \]
\end{formula}

\topic{Partitions of an Integer}
A partition of a positive integer $r$ is a way of writing $r$ as a sum of positive integers, where the order of the summands does not matter. Each summand is called a part. We denote the number of partitions of $r$ by $p(r)$, and the number of partitions of $r$ into exactly $n$ parts by $p(r,n)$.

The numbers $p(r,n)$ satisfy the following conditions:
\begin{formula}
\begin{enumerate}[label=\arabic*.]
\item $p(0,0)=1$, $p(r,0)=0$ for $r>0$, and $p(0,n)=0$ for $n>0$;
\item $p(r,n)=p(r-1,n-1)+p(r-n,n)$ for $r,n\ge 1$.
\end{enumerate}
\end{formula}
\begin{formula}
The number of partitions of $r$ into $n$ parts is the same as the number of partitions of $r$ into parts with maximum size $n$.
\end{formula}

%% ======================= CH 3 =======================
\chapterhead*{Ch 3. Binomial Theorem and Generating Functions}

\topic{Generating Functions}
Given an infinite sequence of numbers $a_0,a_1,a_2,\dots$, we associate with it the (formal) power series
\begin{formula}
\[ F(x)=a_0+a_1x+a_2x^2+\cdots=\sum_{n=0}^{\infty}a_nx^n. \]
\end{formula}
$F(x)$ is called the \key{generating function} of the sequence $a_0,a_1,\dots$. We write $[x^r]F(x)$ for the coefficient of $x^r$ in $F(x)$, so that $a_r=[x^r]F(x)$.
\begin{definition}[3 Common Generating Functions]
\[ 1:\quad A(x)=\sum_{n=0}^{\infty}x^n=\frac{1}{1-x}; \]
\[ n:\quad A(x)=\sum_{n=0}^{\infty}nx^n=\frac{x}{(1-x)^2}; \]
\[ r^n:\quad A(x)=\frac{1}{1-rx}. \]
\end{definition}

\topic{Product Interpretation}
For $i=1,\dots,m$ let $A^{(i)}\subseteq\Zge$ and define its characteristic generating function
\[ F^{(i)}(x)=\sum_{k\in A^{(i)}}x^k. \]
Let $F(x)=F^{(1)}(x)F^{(2)}(x)\cdots F^{(m)}(x)$. Then the coefficient $c_r$ of $x^r$ in $F(x)$ counts the solutions of
\begin{theorem}[Product Interpretation]
\[ x_1+x_2+\cdots+x_m=r,\quad x_i\in A^{(i)}. \]
\end{theorem}
The number of solutions of $x_1+2x_2+5x_5=10$ with $x_1,x_2,x_5\in\Zge$, where $x_1,x_2,x_5$ are the numbers of coins of each kind. By the weighted construction,
\begin{formula}
\[ F(x)=(1+x+x^2+\cdots)(1+x^2+x^4+\cdots)(1+x^5+x^{10}+\cdots) \]
\[ =\frac{1}{(1-x)(1-x^2)(1-x^5)}. \]
\end{formula}
We read off $[x^{10}]F(x)$ by hand, choosing the power from the last factor first and using (AP) over the three cases.

\topic{The Encoding Dictionary}
By the product interpretation, if the $i$-th variable is allowed to take values in $A^{(i)}\subseteq\Zge$, then its factor is $\sum\nolimits_{k\in A^{(i)}}x^k$. The constraints that occur most often, and their factors, are:
\begin{formula}
\begin{cstable}{LCC}
\textbf{constraint on $x_i$} & \textbf{factor} & \textbf{closed form} \\ \midrule
no constraint & $1+x+x^2+\cdots$ & $\frac{1}{1-x}$ \\ \midrule
$x_i\ge a$ & $x^a+x^{a+1}+\cdots$ & $\frac{x^a}{1-x}$ \\ \midrule
$x_i\le b$ & $1+x+\cdots+x^b$ & $\frac{1-x^{b+1}}{1-x}$ \\ \midrule
$a\le x_i\le b$ & $x^a+\cdots+x^b$ & $\frac{x^a-x^{b+1}}{1-x}$ \\ \midrule
$x_i$ even & $1+x^2+x^4+\cdots$ & $\frac{1}{1-x^2}$ \\ \midrule
$x_i$ a multiple of $k$ & $1+x^k+x^{2k}+\cdots$ & $\frac{1}{1-x^k}$
\end{cstable}
\end{formula}
Multiply the factors; the answer is a coefficient of the product.

\topic{Partition Generating Function}
\begin{theorem}[Partition Generating Function]
\[ \sum_{r=0}^{\infty}p(r)x^r=\prod_{k=1}^{\infty}(1+x^k+x^{2k}+\cdots)=\prod_{k=1}^{\infty}\frac{1}{1-x^k}. \]
\end{theorem}

\topic{Exponential Generating Functions}
The exponential generating function (EGF) for the sequence of numbers $(a_r)$ is the power series
\begin{formula}
\[ a_0+a_1\frac{x}{1!}+a_2\frac{x^2}{2!}+a_3\frac{x^3}{3!}+\cdots+a_r\frac{x^r}{r!}+\cdots \]
\[ =\sum_{r=0}^{\infty}a_r\frac{x^r}{r!}. \]
\end{formula}
Thus $a_r=r!\,c_r$, where $c_r$ is the $r$-th coefficient of the power series.

\topic{The EGF Dictionary}
Let $a_r$ denote the number of $r$-permutations of the multiset
\[ \{n_1\cdot b_1,\ n_2\cdot b_2,\ \dots,\ n_k\cdot b_k\}. \]
Then the exponential generating function for $(a_r)$ is
\begin{formula}
\[ \left(1+\frac{x}{1!}+\cdots+\frac{x^{n_1}}{n_1!}\right)\left(1+\frac{x}{1!}+\cdots+\frac{x^{n_2}}{n_2!}\right) \]
\[ \cdots\left(1+\frac{x}{1!}+\cdots+\frac{x^{n_k}}{n_k!}\right), \]
\end{formula}
one factor per type $b_j$, and $a_r$ is the coefficient of $\frac{x^r}{r!}$ in the expansion of the product.
\begin{formula}
\begin{cstable}{>{\raggedright\arraybackslash}p{0.34\linewidth}C}
\textbf{type $b_j$ may appear} & \textbf{factor} \\ \midrule
any number of times & $1+\frac{x}{1!}+\frac{x^2}{2!}+\cdots=e^x$ \\ \midrule
at most $n$ times & $1+\frac{x}{1!}+\cdots+\frac{x^n}{n!}$ \\ \midrule
at least once & $e^x-1$ \\ \midrule
an even number of times & $1+\frac{x^2}{2!}+\frac{x^4}{4!}+\cdots=\frac{e^x+e^{-x}}{2}$ \\ \midrule
an odd number of times & $\frac{x}{1!}+\frac{x^3}{3!}+\cdots=\frac{e^x-e^{-x}}{2}$
\end{cstable}
\end{formula}

%% ======================= CH 4 =======================
\chapterhead*{Ch 4. Recurrence Relations}

\begin{itemize}
\item \key{General method to solve linear recurrence relations.} The general solution is
\begin{formula}
\[ g(n)=g_0(n)+h(n). \]
\end{formula}
\textbf{Finding a particular solution: what to try.} The following table summarises the trial forms for $g_0(n)$. Here $P_k$ and $Q_k$ denote polynomials of degree $k$; the coefficients of $Q_k$ and the constant $C$ are the unknowns. The third column applies when the base $a$ of the exponential (for a polynomial $f$, take $a=1$) is a characteristic root of multiplicity $m$: the trial form must then be multiplied by $n^m$.
\begin{formula}
\begin{cstable}{CCC}
\textbf{$f(n)$} & \textbf{trial $g_0(n)$} & \textbf{if $a$ is a root of multiplicity $m$} \\ \midrule
$P_k(n)$ (here $a=1$) & $Q_k(n)$ & $n^mQ_k(n)$ \\ \midrule
$ca^n$ & $Ca^n$ & $Cn^ma^n$ \\ \midrule
$P_k(n)a^n$ & $Q_k(n)a^n$ & $n^mQ_k(n)a^n$ \\ \midrule
$f_1(n)+f_2(n)$ & $g_{0,1}(n)+g_{0,2}(n)$ & (superposition)
\end{cstable}
\end{formula}

\item \key{Solving linear recurrences using generating functions.} To solve a linear recurrence relation using generating functions:
\begin{enumerate}[label=\arabic*.]
\item Multiply the recurrence by $x^n$ and sum over all $n$ for which it holds. Express each sum in terms of $A(x)$: this gives a linear equation for $A(x)$ (the initial conditions enter here).
\item Solve the linear equation for $A(x)$; it will be a rational function.
\item Extract the coefficient of $x^n$: factor the denominator, use partial fractions, and expand each piece as a power series.
\end{enumerate}
Useful bookkeeping: for the sequence $(a_n)$ with OGF $A(x)$,
\begin{formula}
\[ \sum_{n\ge 1}a_{n-1}x^n=xA(x),\quad \sum_{n\ge 2}a_{n-2}x^n=x^2A(x), \]
\[ \sum_{n\ge k}a_nx^n=A(x)-a_0-a_1x-\cdots-a_{k-1}x^{k-1}. \]
\end{formula}

\item \key{System of recurrence relations.}

\textbf{Method 1:} eliminate $b_n$ to get a single recurrence.

\textbf{Method 2:} generating functions turn the system into linear equations.

Let $A(x)=\sum\nolimits_{n=0}^{\infty}a_nx^n$ and $B(x)=\sum\nolimits_{n=0}^{\infty}b_nx^n$ (with $a_0=1$, $b_0=0$). Multiply each recurrence by $x^n$ and sum over $n\ge 1$:
\begin{formula}
\[ a_n=b_n+a_{n-1}+b_{n-1}+a_{n-2}\Rightarrow A-1=B+xA+xB+x^2A, \]
\[ b_n=a_{n-1}+b_{n-1}\Rightarrow B=xA+xB. \]
\end{formula}
The second gives $B=\frac{x}{1-x}A$; substituting into the first and multiplying by $1-x$,
\begin{formula}
\[ [(1-x-x^2)(1-x)-x(1+x)]A=1-x, \]
\[ A(x)=\frac{1-x}{1-3x-x^2+x^3}, \]
\[ B(x)=\frac{x}{1-x}A(x)=\frac{x}{1-3x-x^2+x^3}. \]
\end{formula}
\end{itemize}

\end{cscolumns}
\end{document}
